Table~\ref{tab:substrate} lists the frozen substrate parameters used
throughout Phase I; the authoritative spec is the V2 protocol
\cite{v2proto,overview}.

\begin{table}[t]
\centering
\caption{\textbf{Substrate specification (frozen V2; the V2 protocol
\cite{v2proto}).} All values commit in the committed configuration
files; later experiments change only the listed flags.}
\label{tab:substrate}
\footnotesize
\begin{tabular}{p{3.4cm}p{5.4cm}p{6.4cm}}
\toprule
Layer / mechanism & Value & Note \\
\midrule
Tick & 1 ms & deterministic; seeded \\
Input channels & 24 (groups \(A=0..7\), \(C=8..15\), \(B=16..23\)) & pure spike sources (D8): no internal state \\
Internal / output neurons & 40 / 12 (ids 24..75, \(N=52\)) & LIF; output is a role tag \\
LIF & \(\tau_m = 20\) ms, \(v_{\mathrm{rest}} = 0\), \(v_{\mathrm{th}} = 1\), refractory 2 ms & \\
Synaptic current & \(\tau_{\mathrm{syn}} = 5\) ms, amplitude \(a = 52\) & current per spike \(= a\cdot w\) \\
Initial wiring & M1: input \(p=0.5\), recurrent \(p=0.2\), weak weights & seeded \\
STDP & pairwise additive; \(\tau_\pm = 20\) ms, \(a^+=0.005\), \(a^-=0.0053\) & E1 amendment A1 operating point \\
Passive decay & \(10^{-6}\)/tick & \\
M2 normalization & per-neuron excitatory budget \(t_e = 0.8\) & multiplicative rescale of live incoming \\
M3 candidates & 6 slots; \(w_{\mathrm{c,init}}=0.01\), \(\theta_{\mathrm{perm}}=0.05\), \(w_{\mathrm{c,perm}}=0.02\) & co-active windows \\
M4 pruning & \(\theta_{\mathrm{prune}} = 0.005\), 10 windows & excitatory only \\
M5 budgets & \(b_e = 40\) exc / \(b_i = 10\) inh per neuron & \\
M6 inhibition & anti-Hebbian; \(w_{\mathrm{inh,max}} = 0.10\) & \\
E6 & per-channel rate balance, history-derived & no labels \\
V2.1 slow state & \(u\): \(+\beta\)/spike, decay \(\tau_s\); \(\beta=0.0046875\), \(\tau_s=5000\) ms & the e24 cell used by all CLLA-era runs \\
Resource detector & runaway: population mean \(>50\) Hz sustained 5 s & abort with failure record \\
Curriculum cadence & 500 ms presentation @ 20 Hz/channel; 1500 ms off & 2 s per presentation \\
Seeds & 20260912, 424242, 9001 & frozen \\
\bottomrule
\end{tabular}
\end{table}


\subsection{Neuron and network dynamics}

Non-input neurons integrate
\begin{equation}
  \frac{dv_i}{dt} = \frac{-(v_i - v_{\mathrm{rest}}) + i_{\mathrm{syn},i} + i_{\mathrm{ext},i}
     - i_{\mathrm{adapt},i} + u_i}{\tau_m},
  \label{eq:dv}
\end{equation}
discretized at 1 ms with exponential kernel decay for
\(i_{\mathrm{syn}}\) (\(\tau_{\mathrm{syn}}\)), adaptation
\(i_{\mathrm{adapt}}\) (E3-era U1; injection
\(\texttt{adaptation\_gain}=0.05\) per spike), and the slow depolarizing
state \(u\) (V2.1),
\begin{equation}
  u_i(t+1) = u_i(t)\, e^{-1/\tau_s} + \beta\cdot
  \mathbf{1}[i \text{ spiked at } t],
\end{equation}
with \(\beta=0\) restoring byte-identical V2 behavior. A neuron spikes
when \(v_i \ge v_{\mathrm{th}}\) after its refractory period; it resets
\(v_i\) to \(v_{\mathrm{rest}}\) and deposits \(a\cdot w\) (excitatory)
or \(-a\cdot w\) (inhibitory) onto each postsynaptic target, with a
one-tick transmission delay. Throughout Phase I the birth triggers were
disabled
(\texttt{birth\_trigger = "none"}); all results concern the fixed
population.

\subsection{Sensory interface and curriculum}

Each input channel produces deterministic Poisson trains (seeded per
pattern, repetition, and channel): 8 channels fire at 20 Hz for 500 ms,
followed by 1500 ms of silence. Patterns \(A\), \(C\), and \(B\) are
channel groups; pattern \(D = A \cup C\) (channels 0--15, same rate) is
used as the composition probe. The curriculum is a configurable stage
schedule; the CLLA-era corpus uses three orders of patterns over three
seeds:
\begin{itemize}
\item \texttt{il}: \(A/C\) interleaved, 20 presentations each (40 total);
\item \texttt{bac}: 20\(\times A\), then 20\(\times C\);
\item \texttt{bca}: 20\(\times C\), then 20\(\times A\);
\item \texttt{d}: 40\(\times D\) (composition probe, informative only
in most CLLA records).
\end{itemize}

\subsection{Plasticity in normative form}

STDP is the pairwise additive trace rule
\cite{bipoo1998,gerstner2002}, applied per tick as
\begin{equation}
  \Delta w = a^+\, \text{pre-trace} \cdot \mathbf{1}[\text{post fired}]
            - a^-\, \text{post-trace} \cdot \mathbf{1}[\text{pre fired}],
\end{equation}
with trailing traces \(\tau_\pm=20\) ms and weights clamped to
\([0,1]\). Where the implemented equations deviate from the normative
form (notably the dimensionality of the allocation gate's drive
estimate; \S\ref{sec:allocation}, \S\ref{sec:repro}), the records are
explicit, and the paper preserves the correction history.